\chapter{Applications and Use Cases}
\label{scenarios}

\section{Primary Use Cases and Applications}

\subsection{Linguistic Field Work}

For endangered languages: automatic pre-annotation provides a first draft.

\begin{enumerate}
\item Run SpeechPrint
\item Review and correct in Praat
\item Export for analysis
\end{enumerate}

\subsection{Subtitle Enhancement for Film and Media}

Video subtitles are typically plain text with no prosodic information. A speaker might ask a question with rising intonation, but the subtitle just shows the words. Using SpeechPrint, one could:

\begin{enumerate}
\item Extract the prosody layer from the audio
\item Add subtle visual markup to the subtitle text (italics for emphasis, punctuation for pauses)
\item Create a richer reading experience that hints at the speaker's tone
\end{enumerate}

This benefits accessibility (helping deaf or hard-of-hearing viewers understand emotional tone) and enables media creators to produce more expressive subtitles.

\subsection{Text-to-Speech with Prosodic Control}

Modern text-to-speech systems (like those in audiobook production) need prosodic cues to sound natural. SpeechPrint can extract the prosody from a human recording of a text, then that prosody profile can be fed to a TTS system as control parameters, making the synthetic speech sound more like the human original.

\subsection{Aesthetic and Artistic Use: Annotated Poetry}

SpeechPrint enables new ways to represent poetry and oral literature. For example:

\begin{enumerate}
\item Record a recitation of a poem
\item Run SpeechPrint to extract phoneme boundaries and prosody
\item Overlay the prosody markup (/, \textbackslash, *, \={}, \_) onto the original poem text
\item Display the annotated poem alongside the audio waveform
\end{enumerate}

This creates a visual representation of how the poem is performed: where emphasis falls, where pitch rises for dramatic effect, where pauses create silence. For literary scholars, this reveals the "sonic topology" of verse (rhythm, rhyme, pitch patterns) that written analysis often misses. For educators, it provides quantitative feedback on expression and control in oral performance.

\section{Conclusions}
\label{ch:conclusions}

\subsection{Summary of Findings}

This thesis presented SpeechPrint, a computational pipeline designed to extract and represent prosodic features from speech recordings in a form accessible to linguistic researchers who lack specialized phonetic training. The system combines modern speech recognition tools (Whisper, WhisperX, Gentle, MFA) with acoustic analysis (Parselmouth) to produce structured TextGrid files containing aligned phonetic and prosodic annotations.

\subsection{Key Contributions}

The main contributions of this work are:

\begin{enumerate}
\item \textbf{Octave error detection and correction}: Systematic evaluation of pitch tracker accuracy, identifying 238 octave errors in Praat's default autocorrelation method on English recordings and demonstrating the superiority of pYIN and CREPE for clean studio speech.

\item \textbf{Integrated annotation pipeline}: A hybrid alignment architecture combining multiple forced aligners (MFA, Gentle, WhisperX) with unified output format, addressing practical limitations of individual tools.

\item \textbf{Symbolic prosody representation}: A minimalist encoding system using ASCII-compatible symbols (\texttt{/}, \texttt{\textbackslash}, \texttt{–}, \texttt{\_}, \texttt{*}) for intuitive visualization of pitch contours and prominence, bridging the gap between phonetic measurement and human interpretation.

\item \textbf{Multilingual evaluation}: Testing across typologically diverse languages (English, German, Daakie, Cabécar) with both high-resource benchmarks (GToBI) and low-resource corpora (DoReCo), demonstrating cross-linguistic applicability.
\end{enumerate}

\subsection{Performance Assessment}

On German GToBI reference sentences, all four pitch trackers achieved agreement within 2\,Hz with no octave errors, confirming that clean studio recordings pose no difficulty for modern methods. English and archival recordings revealed systematic biases in Praat's raw output, successfully corrected by pYIN and CREPE. The final SpeechPrint configuration (pYIN for read speech, CREPE for archival recordings) achieved partial-to-full correctness on benchmark sentences, with consistent pitch direction labeling.

\subsection{Implications for Linguistic Annotation}

SpeechPrint addresses a concrete gap in the linguistic annotation workflow: the absence of prosodic pre-analysis in tools like ELAN. By providing automatic suggestions for phoneme boundaries, syllable structure, and pitch patterns, the system reduces the cognitive load on human annotators while preserving their authority over final annotations. This is particularly valuable for linguistic fieldwork on under-resourced languages, where annotators may lack phonetic training.

\subsection{Objectives Achieved}

The main objective of this thesis was the development of SpeechPrint, a tool that combines phonemic segmentation with phonetic and prosodic analysis. Specific objectives included:

\begin{itemize}
\item Developing a reproducible speech analysis workflow.
\item Combining transcription, forced alignment, and acoustic analysis.
\item Extracting F0, amplitude, duration, and vowel-related information.
\item Testing the system across multiple languages.
\item Investigating prosodic annotation and representation.
\item Exploring artistic, linguistic, and technological use cases.
\item Supporting annotation workflows for linguistic corpora and endangered language documentation.
\end{itemize}

These objectives have been addressed throughout this thesis, with varying degrees of completion and success as detailed in the preceding chapters.

\subsection{Limitations of the Current Implementation}

Despite these contributions, several limitations warrant acknowledgment:

\begin{itemize}
\item Alignment quality remains sensitive to recording conditions; audio with significant background noise or overlapping speakers requires manual correction.
\item Pitch extraction accuracy depends on voice quality; creaky or breathy phonation can produce unreliable F0 estimates even from modern trackers.
\item The symbolic prosody layer encodes only three pitch categories (rising, level, falling), necessarily reducing the granularity of phonetic description compared to systems like ToBI.
\item Cross-linguistic generalization has been tested on a limited sample; performance on languages with fundamentally different prosodic structures (e.g., truly tonal languages) remains unexplored.
\end{itemize}

\section{Future Directions and Open Problems}
\label{ch:discussion}

\subsection{Prosody Resynthesis: From Annotation to Sound}
\label{sec:resynthesis}

In 1769 Wolfgang von Kempelen built the Mechanical Turk, an automaton that
appeared to play chess. Shortly after, he built a second machine: die
Sprech-Maschine, described in his \textit{Mechanismus der menschlichen
Sprache} \cite{kempelen1791}, assembled from bellows, reeds, and a leather
vocal tract shaped to produce vowels and a small number of consonants. The
conjunction of the two inventions implies something precise: without die
Sprech-Maschine preceding it with a voice, the Turk was a clever trick in
a wooden box. The voice was not an ornament on the rational machine. It was
what made the rational machine believable. Mladen Dolar, writing about the
metaphysics of voice, calls this ``the most human of effects, an effect of
`interiority'\,'' \cite[10]{dolar2006}: the sense that something is present
behind the mechanism, that a subject inhabits the sound. Two and a half
centuries later, the same question returns. Speech recognition and synthesis
systems now transcribe and generate language at a level indistinguishable
from human production in many conditions. What they do not yet recover or
produce reliably is the prosodic layer: the emphasis, the question, the
hesitation, the attitude. The voice without prosody is still, in some sense,
the Turk without the Sprech-Maschine.

A natural extension of the SpeechPrint pipeline is to close the loop between
analysis and synthesis. The current system moves in one direction: from audio
to symbolic annotation. The inverse direction (from symbolic annotation to
audio) would constitute a \emph{prosody resynthesis} capability. This concept
is illustrated in Figure~\ref{fig:resynthesis_pipeline}.

The proposed pipeline would work as follows. First, a recording is processed by
the existing SpeechPrint pipeline to produce a TextGrid with a prosody tier.
Second, a user modifies the prosody symbols in the TextGrid, for example,
changing a falling accent (\texttt{*\textbackslash}) to a rising one
(\texttt{*/}) on the word \textit{Mary} to produce the reading ``MARY flew to
Milan yesterday'' rather than ``Mary FLEW to Milan yesterday.'' Third, the
modified symbolic annotation is translated into prosody control parameters that
are passed to a speech synthesis system. The synthesis system produces a new
audio file with the requested prosodic pattern while preserving the voice quality
and segmental content of the original.

The linguistic term for this operation is \emph{prosody transfer} or
\emph{intonation resynthesis}. It has been studied in the context of
human-to-human style transfer and as a component of expressive text-to-speech
synthesis \cite{skerry2018towards}. The key challenge is the mapping from
symbolic prosody (the / and \textbackslash\ labels of SpeechPrint's tier) to
quantitative F0 targets that a synthesis system can use. One approach is to
learn this mapping from a corpus of annotated speech; another is to use a
parametric F0 model in which the symbolic labels directly control the slope,
height, and velocity of F0 movements.

A practical implementation using currently available tools would use the
Coqui TTS framework or a similar system that exposes duration and pitch control
parameters. Coqui XTTS, for example, allows voice cloning from a short reference
segment and accepts prosodic modifications through its Python API. The SpeechPrint
symbolic tier would be compiled into a sequence of pitch targets at word or
syllable boundaries, which the synthesis system would interpolate into a smooth
F0 contour. This would allow a researcher to hear the prosodic consequences of an
annotation change without recording a new utterance, a capability that could be
valuable both for validation of the annotation and for generating stimuli in
prosody experiments.

\begin{figure}[htbp]
\centering
\begin{tabular}{|c|}
\hline
\textbf{Analysis Direction (Current)} \\[0.5em]
Audio File $\rightarrow$ Speech Recognition (WhisperX) $\rightarrow$ \\
Forced Alignment (MFA / Gentle) $\rightarrow$ Pitch Extraction \\
(pYIN or CREPE) $\rightarrow$ Symbolic Prosody (/, \textbackslash, *, \_) $\rightarrow$ \\
Annotated TextGrid \\[1em]
\hline
\textbf{Resynthesis Direction (Proposed)} \\[0.5em]
Edited TextGrid $\rightarrow$ F0 Target Compiler $\rightarrow$ \\
Speech Synthesis (Coqui TTS/XTTS) $\rightarrow$ Resynthesized Audio \\
\hline
\end{tabular}
\caption{The proposed bidirectional SpeechPrint pipeline. The analysis direction
is fully implemented and demonstrated on multiple languages. The resynthesis direction
(future work) would enable users to modify prosody annotations and hear the perceptual
consequences without re-recording.}
\label{fig:resynthesis_pipeline}
\end{figure}

A recent system called DrawSpeech \cite{chen2025drawspeech} provides
direct empirical support for this direction. DrawSpeech conditions a
latent diffusion model on prosodic sketches drawn by users: instead of
specifying a full F0 contour, the user draws a rough line indicating the
desired pitch trend per word, and the model recovers the detailed contour
from that coarse input. Experimental results show that sketch-conditioned
synthesis achieves a mean opinion score of 4.49 (slightly above the
4.46 ground-truth reference) while also enabling precise per-word
prosody control. Their Figure~3 demonstrates four different sketches
applied to the sentence ``I didn't say you stole the money'', each
emphasising a different word, producing perceptually distinct utterances.
This is precisely the minimal-pairs corpus used in this thesis, but as
synthesised output rather than recorded speech.

The symbolic prosody tier
produced by SpeechPrint (per-syllable labels such as
\texttt{*{\={}}/}, accented, high, rising) functions as a
human-readable version of the DrawSpeech sketch. The analysis pipeline
described in this thesis extracts the sketch from recorded speech; a
DrawSpeech-style synthesis system would use an edited version of that
sketch to generate new speech. The two systems are complementary halves
of a closed loop: analyse prosody from audio, edit it symbolically, and
synthesise it back.

\subsection{A Speech DAW: Interactive Prosody Manipulation}
\label{sec:speech_daw}

Taking the resynthesis concept further, a \emph{speech digital audio workstation}
(Speech DAW) would present the TextGrid annotation tiers as an editable timeline,
similar in concept to how a music DAW presents MIDI notes. The user would be able
to click on any syllable label in the prosody tier and change it, drag the onset
or offset of a pitch movement, raise or lower the height symbol from \texttt{\_}
to \texttt{{\={}}}, or add or remove an accent marker \texttt{*}. Each edit would
trigger a re-synthesis of the affected region, allowing real-time auditory
feedback.

The design is motivated by a gap in current tools. Praat allows the manual
manipulation of a pitch track at the acoustic level, which requires expert
knowledge of what a changed F0 contour sounds like. The proposed Speech DAW
would operate at the symbolic level, where the units (/ \textbackslash\, *) are
interpretable by a non-specialist. This lowers the barrier for users in
conversation analysis or field linguistics who want to test hypotheses about
prosodic structure by listening to counterfactual examples.

\subsection{Automatic Prosody Classification}
\label{sec:prosody_classifier}

A third future direction is the development of a classifier that assigns one of a
small number of prosodic categories to each sentence or utterance. The motivation
is that a single global label (whether a sentence is a statement, a
question, an emphatic assertion, a continuation, or a boundary) is often more
useful for annotation workflows than a fine-grained per-syllable sequence of
symbols. The classification would be trained on existing annotated corpora and
applied to new recordings.

A candidate set of five prosodic categories, sufficient to cover the most
frequent intonational events in English and German, might be:
\begin{enumerate}
\item \textbf{Declarative}: broad-focus statement with final fall.
\item \textbf{Interrogative}: yes/no question with final rise.
\item \textbf{Narrow focus}: prominent accent on a single content word,
      marking contrast or correction.
\item \textbf{Continuation}: non-final intonational phrase with
      non-falling boundary tone, signalling that more speech follows.
\item \textbf{Exclamation}: emphatic or expressive contour, typically with
      high onset and rapid fall.
\end{enumerate}

This classification scheme is a simplification of the GToBI inventory but
captures the categories most relevant to the use cases described in
Chapter~\ref{scenarios}. The classifier would take the sequence of SpeechPrint
symbolic labels as input features, supplemented by the mean F0, F0 range, and
final boundary tone of the utterance. The output category would be displayed as
an additional tier in the TextGrid, providing a sentence-level interpretive
summary alongside the syllable-level labels.

\subsection{Bias, Projection, and Cross-Lingual Phonological Approximation}
\label{sec:bias_projection}

A recurring observation in this thesis is that applying a speech model trained
on one language to a recording in an unrelated language does not produce random
noise: it produces something coherent but wrong. When Whisper's Italian
acoustic model was applied to the Daakie recording, it returned phonetically
plausible Italian word sequences, not silence or unintelligible garbage.
The model could not understand Daakie, but it
projected its learned phonological structure onto the signal and produced the
nearest thing in its training distribution. This is model bias operating as a
generative mechanism: the failure is systematic and, viewed from a different
angle, compositionally interesting. Each reinterpretation layer introduces a
transformation that is determined by what the model knows and what it does not,
producing emergent changes in apparent timbre, prosody, and intelligibility that
are not random but structured by the model's prior. The SpeechPrint prosody
layer applied to such a mistranscribed signal assigns labels that are acoustically
correct but linguistically misaligned, creating a kind of prosodic ghost of the
original utterance. This phenomenon (the gap between what a model hears and
what was said) represents a potential compositional space for work at the
intersection of speech technology and sound art, and connects to a broader
research thread on how bias and error in computational listening shape the
transformation of voice across systems.
